\paragraph{临界五次方程的伽罗瓦障碍与半平面根分布}
在定理 \ref{thm:xi-rh-critical-before-nearest-complex-branch} 与注 \ref{rem:collision-cubic-rhsharp-quintic} 的阈值多项式
\[
P(u):=u^5+4u^4+3u^3-96u^2-42u-14
\]
之上，还可进一步把唯一正实临界根 \(u_R\) 的代数复杂度与其余四根的稳定半平面位置完全锁定。

\begin{proposition}[临界五次方程的伽罗瓦障碍与 \texorpdfstring{$u_R$}{uR} 的根式不可解性]\label{prop:xi-collision-threshold-quintic-unsolvability}
\leanverified{paper\_xi\_collision\_threshold\_quintic\_unsolvability}
设 \(u_R>1\) 为定理 \ref{thm:xi-rh-critical-before-nearest-complex-branch} 中满足 \(P(u_R)=0\) 的唯一正实根。
则 \(P\) 在 \(\QQ\) 上不可约，且其伽罗瓦群满足
\[
\Gal(P/\QQ)\cong S_5.
\]
因此 \(u_R\) 不能由有理数经有限次四则运算与开方、开立方等根式运算表达。
\end{proposition}
\begin{proof}
先看不可约性。把 \(P\) 模 \(11\) 约化，直接检验可知其在 \(\FF_{11}[u]\) 中既无一次因子也无二次因子。
由于 \(\deg P=5\)，故 \(P\bmod 11\) 必为不可约五次，进而由 Gauss 引理知 \(P\) 在 \(\QQ\) 上不可约。

再看伽罗瓦群。模 \(11\) 不可约说明
\(
G:=\Gal(P/\QQ)\subset S_5
\)
是传递子群且含有一个 \(5\)-循环。
另一方面，模 \(31\) 约化时 \(P\) 分解为一个不可约二次因子与一个不可约三次因子之积。
由 Dedekind 分解型判据，\(G\) 含有循环型为 \((2)(3)\) 的元素；其平方给出一个 \(3\)-循环。

最后计算判别式：
\[
\mathrm{Disc}(P)=2^6\cdot 3^2\cdot 7\cdot 743\cdot 1693^2,
\]
它不是 \(\QQ\) 中的平方，故 \(G\nsubseteq A_5\)。
而 \(S_5\) 的含 \(5\)-循环的传递子群只有
\[
C_5,\qquad D_5,\qquad F_{20},\qquad A_5,\qquad S_5.
\]
其中 \(C_5,D_5,F_{20}\) 均不含 \(3\)-循环，\(A_5\) 又已由判别式非平方排除，
故唯一可能为 \(G=S_5\)。

由 Abel--Ruffini 定理，\(P\) 的根不可能由根式统一表达；特别地，其唯一正实根 \(u_R\) 亦不可由根式表达。证毕。
\end{proof}

\begin{proposition}[临界五次方程的唯一右半平面根与左半平面余根]\label{prop:xi-collision-threshold-quintic-halfplane-root-distribution}
\leanverified{paper\_xi\_collision\_threshold\_quintic\_halfplane\_root\_distribution}
对同一多项式
\[
P(u)=u^5+4u^4+3u^3-96u^2-42u-14,
\]
开右半平面 \(\{\Re(u)>0\}\) 内恰有一个零点；该零点正是 \(u_R\)。
其余四个零点全部满足
\[
\Re(u)<0.
\]
\end{proposition}
\begin{proof}
对 \(P\) 作 Routh 表，其首列为
\[
1,\qquad
4,\qquad
27,\qquad
-\frac{2438}{27},\qquad
-\frac{104069}{2438},\qquad
-14.
\]
该首列无零项且仅发生一次变号。
因此由 Routh--Hurwitz 判据，\(P\) 在开右半平面内恰有一个零点，并且不存在纯虚根。

另一方面，\(P\) 的系数符号序列为
\[
(+,+,+,-,-,-),
\]
仅有一次变号。
由 Descartes 符号法则，\(P\) 恰有一个正实根。
结合定理 \ref{thm:xi-rh-critical-before-nearest-complex-branch} 中对临界参数的唯一性，这个正实根正是 \(u_R\)。

由于 \(P\) 的系数为实数，非实根必成共轭对出现。
既然右半平面内只有一个零点且不存在纯虚根，除 \(u_R\) 之外的其余四根便只能全部落在开左半平面。证毕。
\end{proof}

\endinput
